%==============
\begin{center}
{\LARGE \bf New Applications in {\em Mathematica}}
\end{center}
%==============
\bigskip

\noindent
{\bf Organizer}
\medskip

A. Bocharov (alexei@wolfram.com)

\bigskip

\noindent
{\bf Description}
\medskip

The session covers new {\em Mathematica} packages and algorithms developed with Mathematica in
mind. Tools for Abstract and Complex Calculus , Symbolic/Numeric methods, tools for algorithm
generation, and some applications to Mechanics are presented. Along with specific applications to
problems in Algebra, Differential Equations, Numerical Analysis, some new Mathematica
features will be shown (although on "as needed" basis). 

\bigskip
\noindent
{\bf Talks}
\medskip

Date: July 20th (Saturday) 

\begin{description}
\item[08:30-09:00] \ \\
  R. Barrere\\
  {\bf Functional Generic Algorithms for Symbolic Approximations.} 
\item[9:00-09:30] \ \\
  A. V. Bocharov, E. M. Vorob'ev\\
  {\bf Variations on the Theme of Thomas-Rosenfeld with the New
    Version of {\em Mathematica}.} 
\item[09:30-10:00] \ \\
  J. Korelc\\
  {\bf A Mathematica Package for Cooperative Problem Solving in Computational
    Mechanics.} 
\item[10:00-10:30] \ \\
  \\
  {\bf Coffee Break} 
\item[10:30-11:10 ] \ \\
  M. Sofroniou\\
  {\bf Computer derivation of numerical methods.} 
\item[11:10-11:40] \ \\
  M. Trott\\
  {\bf Computation and visualization of general Riemann surfaces w[z] of arbitrary algebraic
    functions.} 
\item[11:40-12:00] \ \\
  Jason Harris\\
  {\bf Extensible Notations in Mathematica 3.0.} 
\end{description}

